% GENERATED FILE. Edit canonical lessons, then run npm run generate:books.
% canonical-source-hash: fnv1a64:8f7008314b208876
% canonical-lessons: TE-C268-pradarsana, TE-C268-rediyo, TE-C268-karyakramam, TE-C268-intarnet, TE-C268-vebsait

\chapter{Show, Radio, Programme, Internet, Website}
\label{ch:te-show-radio-programme-internet-website}
\hlchaptermodality{268}

\begin{chapteropening}
\textbf{By the end of this chapter:} \emph{I can say a show, the radio, a programme, the internet and a website.}

Complete the last lesson of chapter 268: I can say a show, the radio, a programme, the internet and a website.
\end{chapteropening}

\section[\texorpdfstring{\te{ప్రదర్శన} (pradarśana)}{ప్రదర్శన (pradarśana)}]{\te{ప్రదర్శన} (pradarśana) — a show, an exhibition}

\label{lesson:TE-C268-pradarsana}



\begin{quote}
Before the new one: say the Telugu for moonlight, then the Telugu for darkness.
\end{quote}

\subsection*{You'll want to know: \te{ప్రదర్శన}}
\textbf{\te{ప్రదర్శన}} — \emph{pradarśana} — \textquotedblleft{}a show, an exhibition\textquotedblright{}.

\begin{grammarlens}[title={What you've built}]
One more word for this chapter.
\end{grammarlens}

\subsection*{Your turn}
\begin{itemize}
\raggedright
  \item \emph{Say it:} \emph{pradarśana}
  \item \emph{Say it:} \emph{pradarśana}, once more
  \item \emph{Say it:} say \emph{pradarśana}
  \item \emph{Recall:} say the Telugu for moonlight, then the Telugu for darkness, then say \emph{pradarśana} again
\end{itemize}

\begin{tcolorbox}[breakable,colback=teal!4,colframe=teal!35!black,title={Before you move on}]
What does \te{ప్రదర్శన} mean? (\textquotedblleft{}A show, an exhibition\textquotedblright{}.) Say it once more.
\end{tcolorbox}
\section[\texorpdfstring{\te{రేడియో} (rēḍiyō)}{రేడియో (rēḍiyō)}]{\te{రేడియో} (rēḍiyō) — the radio}

\label{lesson:TE-C268-rediyo}



\begin{quote}
Before the new one: say the Telugu for darkness, then the Telugu for a show.
\end{quote}

\subsection*{You'll want to know: \te{రేడియో}}
\textbf{\te{రేడియో}} — \emph{rēḍiyō} — \textquotedblleft{}the radio\textquotedblright{}.

\begin{grammarlens}[title={What you've built}]
One more word for this chapter.
\end{grammarlens}

\subsection*{Your turn}
\begin{itemize}
\raggedright
  \item \emph{Say it:} \emph{rēḍiyō}
  \item \emph{Say it:} \emph{rēḍiyō}, once more
  \item \emph{Say it:} say \emph{rēḍiyō}
  \item \emph{Recall:} say the Telugu for darkness, then the Telugu for a show, then say \emph{rēḍiyō} again
\end{itemize}

\begin{tcolorbox}[breakable,colback=teal!4,colframe=teal!35!black,title={Before you move on}]
What does \te{రేడియో} mean? (\textquotedblleft{}The radio\textquotedblright{}.) Say it once more.
\end{tcolorbox}
\section[\texorpdfstring{\te{కార్యక్రమం} (kāryakramaṁ)}{కార్యక్రమం (kāryakramaṁ)}]{\te{కార్యక్రమం} (kāryakramaṁ) — a programme, a show}

\label{lesson:TE-C268-karyakramam}



\begin{quote}
Before the new one: say the Telugu for a show, then the Telugu for the radio.
\end{quote}

\subsection*{You'll want to know: \te{కార్యక్రమం}}
\textbf{\te{కార్యక్రమం}} — \emph{kāryakramaṁ} — \textquotedblleft{}a programme, a show\textquotedblright{}.

\begin{grammarlens}[title={What you've built}]
One more word for this chapter.
\end{grammarlens}

\subsection*{Your turn}
\begin{itemize}
\raggedright
  \item \emph{Say it:} \emph{kāryakramaṁ}
  \item \emph{Say it:} \emph{kāryakramaṁ}, once more
  \item \emph{Say it:} say \emph{kāryakramaṁ}
  \item \emph{Recall:} say the Telugu for a show, then the Telugu for the radio, then say \emph{kāryakramaṁ} again
\end{itemize}

\begin{tcolorbox}[breakable,colback=teal!4,colframe=teal!35!black,title={Before you move on}]
What does \te{కార్యక్రమం} mean? (\textquotedblleft{}A programme, a show\textquotedblright{}.) Say it once more.
\end{tcolorbox}
\section[\texorpdfstring{\te{ఇంటర్నెట్} (iṇṭarneṭ)}{ఇంటర్నెట్ (iṇṭarneṭ)}]{\te{ఇంటర్నెట్} (iṇṭarneṭ) — the internet}

\label{lesson:TE-C268-intarnet}



\begin{quote}
Before the new one: say the Telugu for the radio, then the Telugu for a programme.
\end{quote}

\subsection*{You'll want to know: \te{ఇంటర్నెట్}}
\textbf{\te{ఇంటర్నెట్}} — \emph{iṇṭarneṭ} — \textquotedblleft{}the internet\textquotedblright{}.

\begin{grammarlens}[title={What you've built}]
One more word for this chapter.
\end{grammarlens}

\subsection*{Your turn}
\begin{itemize}
\raggedright
  \item \emph{Say it:} \emph{iṇṭarneṭ}
  \item \emph{Say it:} \emph{iṇṭarneṭ}, once more
  \item \emph{Say it:} say \emph{iṇṭarneṭ}
  \item \emph{Recall:} say the Telugu for the radio, then the Telugu for a programme, then say \emph{iṇṭarneṭ} again
\end{itemize}

\begin{tcolorbox}[breakable,colback=teal!4,colframe=teal!35!black,title={Before you move on}]
What does \te{ఇంటర్నెట్} mean? (\textquotedblleft{}The internet\textquotedblright{}.) Say it once more.
\end{tcolorbox}
\section[\texorpdfstring{\te{వెబ్సైట్} (vebsaiṭ)}{వెబ్సైట్ (vebsaiṭ)}]{\te{వెబ్సైట్} (vebsaiṭ) — a website}

\label{lesson:TE-C268-vebsait}



\begin{quote}
Before the new one: say the Telugu for a programme, then the Telugu for the internet.
\end{quote}

\subsection*{You'll want to know: \te{వెబ్సైట్}}
\textbf{\te{వెబ్సైట్}} — \emph{vebsaiṭ} — \textquotedblleft{}a website\textquotedblright{}.

\begin{grammarlens}[title={What you've built}]
One more word for this chapter.
\end{grammarlens}

\subsection*{Your turn}
\begin{itemize}
\raggedright
  \item \emph{Say it:} \emph{vebsaiṭ}
  \item \emph{Say it:} \emph{vebsaiṭ}, once more
  \item \emph{Say it:} say \emph{vebsaiṭ}
  \item \emph{Recall:} say the Telugu for a programme, then the Telugu for the internet, then say \emph{vebsaiṭ} again
\end{itemize}

\begin{tcolorbox}[breakable,colback=teal!4,colframe=teal!35!black,title={Before you move on}]
What does \te{వెబ్సైట్} mean? (\textquotedblleft{}A website\textquotedblright{}.) Say it once more.
\end{tcolorbox}
